\origpage{40}{ROBERTSON AND SEYMOUR}
There are three attractive features of the theorem which serve as its
motivation. First, the theorem is, in a sense, sharp, for it can be reformulated
as follows:
\stmt
(1.2). \textit{Let $\mathscr{F}$ be a set of graphs. Then the following are equivalent}:

\hspace*{12pt}(i)\enspace \textit{for some integer $w$, every $G\in\mathscr{F}$ has path-width $\le w$};

\hspace*{12pt}(ii)\enspace \textit{there is a forest $H$ such that no $G\in\mathscr{F}$ has a minor isomorphic
to $H$.}

It is easy to prove the equivalence of (1.2) with (1.1) by making use of
two observations: (a) if $H$ is a minor of $G$, then its path-width is no greater
than the path-width of $G$, and (b) there are trees with arbitrarily large path-width
(e.g., the tree $\mathbf{Y}_\lambda$ defined later has path-width $\lceil\frac12(\lambda+1)\rceil$, for any
integer $\lambda\ge 1$, where $\lceil x\rceil$ is the least integer not less than $x$).

\fix{We give the details. For (a), let $H$ be isomorphic to a minor of $G$. Then $H$ is isomorphic
to a graph obtained from a subgraph $G'$ of $G$ by contracting a set $C$ of edges. For each
$v\in V(H)$ let $Y(v)$ be the vertex set of the component of the graph with vertex set $V(G')$ and
edge set $C$ that is contracted to $v$, and for each edge $e$ of $H$ let $f(e)$ be the
corresponding edge of $G'$. Then the sets $Y(v)$ $(v\in V(H))$ are disjoint, each $G|Y(v)$ is
connected, and for each edge $e$ of $H$ with ends $v,v'$, the edge $f(e)$ has one end in $Y(v)$
and the other in $Y(v')$. (These are the conditions (M1)--(M3) of Section 3.) We use the
following.

\hspace*{12pt}(1) \textit{Let $X_1,\dots,X_r$ be a path-decomposition of $G$, let $K$ be a
connected subgraph of $G$, and let $1\le i\le j\le l\le r$. If $K$ has a vertex in $X_i$ and a
vertex in $X_l$, then $K$ has a vertex in $X_j$.}

Suppose not, and let $x\in V(K)\cap X_i$ and $y\in V(K)\cap X_l$. We claim that every vertex of
$K$ is in some $X_t$ with $t\le j$. This holds for $x$, since $i\le j$. Let $z,w$ be adjacent
vertices of $K$, where $z\in X_t$ for some $t\le j$. By (W1) some $X_s$ contains both $z$ and $w$.
If $s>j$, then $z\in X_t\cap X_s\subseteq X_j$ by (W2), which is impossible since $K$ has no vertex
in $X_j$. So $s\le j$, and $w\in X_s$. Since $K$ is connected, every vertex of $K$ is reached from
$x$ by such steps, and the claim follows. In particular $y\in X_t$ for some $t\le j$, and so
$y\in X_t\cap X_l\subseteq X_j$ by (W2), which is impossible. This proves (1).

Now let $X_1,\dots,X_r$ be a path-decomposition of $G$ with $|X_i|\le k+1$ $(1\le i\le r)$. For
$1\le i\le r$ let $X'_i$ be the set of all $v\in V(H)$ with $Y(v)\cap X_i\ne\varnothing$. We claim
that $X'_1,\dots,X'_r$ is a path-decomposition of $H$ with $|X'_i|\le k+1$ $(1\le i\le r)$. For
(W1), let $e$ be an edge of $H$ with ends $v,v'$. Then $f(e)$ has an end $a\in Y(v)$ and an end
$b\in Y(v')$, and by (W1) for $G$ some $X_i$ contains $a$ and $b$. So $v,v'\in X'_i$. For (W2), let
$1\le i\le j\le l\le r$ and $v\in X'_i\cap X'_l$. Then $G|Y(v)$ is a connected subgraph of $G$ with
a vertex in $X_i$ and a vertex in $X_l$, and so it has a vertex in $X_j$, by (1). Thus $v\in X'_j$.
For the sizes, choose for each $v\in X'_i$ a vertex in $Y(v)\cap X_i$. Since the sets $Y(v)$ are
disjoint, distinct members of $X'_i$ give distinct vertices of $X_i$, and so
$|X'_i|\le|X_i|\le k+1$. This proves (a).

For (b) we prove that $\mathbf{Y}_{2k+1}$ has path-width $>k$ for every integer $k\ge 0$. (This is
the lower bound above for odd $\lambda$, which is all that we need.) For $k=0$ it is true, since
$\mathbf{Y}_1$ has an edge, and by (W1) some set of any path-decomposition contains both of its
ends. Let $k\ge 0$, suppose that $\mathbf{Y}_{2k+1}$ has path-width $>k$, and let
$\mathbf{Y}=\mathbf{Y}_{2k+3}$. Let $c$ be the vertex of valency 3 of the copy of $K_{1,3}$ from
which $\mathbf{Y}$ is built, and let $K_1, K_2, K_3$ be the components of $\mathbf{Y}\backslash\{c\}$.
Each $K_j$ is isomorphic to the tree $\mathbf{H}_{2k+2}$ of Section 3, and so it has a minor
isomorphic to $\mathbf{Y}_{2k+1}$, by (3.3). Hence $K_j$ has path-width $>k$, by (a). Now suppose
that $X_1,\dots,X_r$ is a path-decomposition of $\mathbf{Y}$ with $|X_i|\le k+2$ $(1\le i\le r)$. For
$j=1,2,3$ the sets $X_i\cap V(K_j)$ $(1\le i\le r)$ form a path-decomposition of $K_j$, and so
$|X_{i_j}\cap V(K_j)|\ge k+2$ for some $i_j$. Then $X_{i_j}\subseteq V(K_j)$ and
$X_{i_j}\ne\varnothing$. Since $K_1,K_2,K_3$ are disjoint, $i_1, i_2, i_3$ are distinct, and we may
assume that $i_1<i_2<i_3$. The graph $\mathbf{Y}\backslash V(K_2)$ is connected, since $c$ is
adjacent to a vertex of $K_1$ and to a vertex of $K_3$. It contains $K_1$ and $K_3$, and so it has
a vertex in $X_{i_1}$ and a vertex in $X_{i_3}$. By (1) it has a vertex in $X_{i_2}$. But
$X_{i_2}\subseteq V(K_2)$, which is impossible. Thus $\mathbf{Y}_{2k+3}$ has path-width $>k+1$, and
(b) follows by induction.

Now assume (1.1), and let $\mathscr{F}$ be a set of graphs. If (i) holds, let
$H=\mathbf{Y}_{2w+1}$, which is a tree. If some $G\in\mathscr{F}$ had a minor isomorphic to $H$,
then $H$ would have path-width $\le w$ by (a), contrary to (b). So (ii) holds. If (ii) holds, (1.1)
applied to $H$ gives (i).\fixnote{fix:1.2-proof-added}{The paper leaves the derivation of (1.2)
from (1.1), and the proofs of the observations (a) and (b), to the reader. The argument above is
the one used in the formalization. It uses (3.3), whose proof does not depend on (1.2).}}

The second attractive feature of (1.1) is that it dovetails nicely with the
main theorem of [3]; the two results together yield the following:
\stmt
(1.3). \textit{Let $G_1, G_2,...,$ be a countable sequence of graphs, such that $G_1$ is a
forest. Then there exist $j>i\ge 1$ such that $G_i$ is isomorphic to a minor of $G_j$.}
\stmt
(This implies the final result mentioned in the abstract, as may easily be
seen.) Theorem (1.3) is relevant to an interesting conjecture of Wagner
(unpublished), that (1.3) is true even without the hypothesis that $G_1$ is a
forest. See [3] for a more complete discussion.

And third (see [1]), (1.1) yields a ``good'' algorithm to test if an arbitrary
graph $G$ has a minor isomorphic to a fixed forest $H$. (``Good'' here is in its
technical sense of ``polynomially bounded.'' It is not a practical algorithm;
the exponent of the polynomial, although constant, is enormous.)

For the purposes of this paper we need the following three facts about
path-width (the proofs are easy).
\stmt
(1.4). \textit{If every connected component of $G$ has path-width $\le k$, then $G$ has
path-width $\le k$.}
\stmt
(1.5). \textit{If $X\subseteq V(G)$ and $G\backslash X$ has path-width $\le k$, then $G$ has path-width $\le k+|X|$.}

($G\backslash X$ denotes the result of deleting $X$.)
\stmt
(1.6). \textit{If $X_1,..., X_r\subseteq V(G)$ is a path-decomposition of $G$, and for
$1\le i\le r-1$, $|X_i\cap X_{i+1}|\le k$, and for $1\le i\le r$,}
\begin{equation*}
G\backslash\bigcup(X_j\colon 1\le j\le r, j\ne i)
\end{equation*}
\textit{has path-width $\le k'$, then $G$ has path-width $\le k'+2k$.}

\fix{We give the proofs. First, if $G$ has path-width $\le k$, then it has a path-decomposition
$X_1,\dots,X_r$ with $|X_i|\le k+1$ $(1\le i\le r)$ in which every vertex is in some $X_i$. For by
(W1) a vertex in no $X_i$ is incident with no edge, and we may add the set $\{v\}$ at the end of
the sequence for each such vertex $v$. Then (W1) and (W2) still hold.

For (1.4), take a path-decomposition of each component of $G$ with sets of size $\le k+1$, and
write these sequences one after another. The result is a path-decomposition of $G$. For every edge
of $G$ is an edge of some component, and each vertex is in sets of only one of the sequences,
since the sets of each sequence are subsets of the vertex set of its component.

For (1.5), take a path-decomposition $X_1,\dots,X_r$ of $G\backslash X$ with $|X_i|\le k+1$
$(1\le i\le r)$ in which every vertex of $G\backslash X$ is in some $X_i$. We claim that
$X_1\cup X,\dots,X_r\cup X,X$ is a path-decomposition of $G$. Its sets have size $\le k+1+|X|$,
and so this will prove (1.5). For (W1), every set contains $X$, and so an edge with both ends in
$X$ has both ends in the last set. An edge with just one end in $X$ has its other end in some
$X_i$, and so both ends in $X_i\cup X$. An edge with no end in $X$ is an edge of $G\backslash X$, and
so both its ends are in some $X_i$. For (W2), a vertex in $X$ is in every set, and a vertex not in
$X$ is in $X_i\cup X$ exactly when it is in $X_i$, and is not in the last set.

For (1.6), let $X_1,\dots,X_r$ be as in (1.6). For $1\le i\le r$ let
$B_i=X_i\cap\bigcup(X_j\colon j\ne i)$ and $G_i=G\backslash\bigcup(X_j\colon j\ne i)$. A vertex in
$X_i\cap X_j$ with $j<i$ is in $X_{i-1}$, and one with $j>i$ is in $X_{i+1}$, by (W2). So
$B_i\subseteq(X_{i-1}\cap X_i)\cup(X_i\cap X_{i+1})$ (where $X_0=X_{r+1}=\varnothing$), and
$|B_i|\le 2k$. By hypothesis and the first paragraph, $G_i$ has a path-decomposition
$Z_{i1},\dots,Z_{im_i}$ with $|Z_{ij}|\le k'+1$ $(1\le j\le m_i)$ in which every vertex of $G_i$ is
in some set. We call the sequence
\begin{equation*}
(Z_{i1}\cap X_i)\cup B_i,\;\dots,\;(Z_{im_i}\cap X_i)\cup B_i,\;B_i
\end{equation*}
the \textit{refinement} of $X_i$. Its sets are subsets of $X_i$ of size $\le k'+1+2k$, and each of
them contains $B_i$.

\hspace*{12pt}(2) \textit{For $1\le i\le r$, the refinement of $X_i$ is a path-decomposition of
$G|X_i$, and every vertex of $X_i$ is in some set of it.}

A vertex of $X_i-B_i$ is in no $X_j$ with $j\ne i$, and so it is a vertex of $G_i$, and it is in
some $Z_{ij}$. The vertices of $B_i$ are in the last set. For (W1), let $e$ be an edge of $G|X_i$
with ends $u,v$. If $u,v\notin B_i$, then $u,v\notin\bigcup(X_j\colon j\ne i)$, so $e$ is an edge of
$G_i$, and some $Z_{ij}$ contains $u$ and $v$, by (W1) for $G_i$. Then $(Z_{ij}\cap X_i)\cup B_i$
contains both. If $u\in B_i$, then $v$ is in some set of the refinement, as we just saw, and $u$
is in every set. The case $v\in B_i$ is the same. For (W2), a vertex of $B_i$ is in every set, and a
vertex $v\in X_i-B_i$ is in $(Z_{ij}\cap X_i)\cup B_i$ exactly when $v\in Z_{ij}$, and is not in the
last set. So (W2) follows from (W2) for $Z_{i1},\dots,Z_{im_i}$. This proves (2).

Now write the refinements of $X_1,\dots,X_r$ one after another. We claim that the result is a
path-decomposition of $G$. Since its sets have size $\le k'+1+2k$, this will prove (1.6). For
(W1), let $e$ be an edge with ends $u,v$. By (W1), $u,v\in X_i$ for some $i$. So $e$ is an edge of
$G|X_i$, and some set of the refinement of $X_i$ contains $u$ and $v$, by (2). For (W2), let $v$ be
in a set of the refinement of $X_i$ and in a later set of the refinement of $X_l$, where $i\le l$,
and consider a set between them, in the refinement of $X_j$ with $i\le j\le l$. If $i=l$, then
$j=i$, and $v$ is in the set by (2). Suppose that $i<l$. The sets of the refinements of $X_i$ and
$X_l$ are subsets of $X_i$ and $X_l$, and so $v\in X_i\cap X_l$, and $v\in X_j$ by (W2). Since $i<l$,
one of $i,l$ is not $j$, and so $v$ is in $X_j$ and in some other set $X_{j'}$. Hence $v\in B_j$. So $v$ is in every
set of the refinement of $X_j$.\fixnote{fix:1.4-1.6-proof-added}{The paper says that the proofs of
(1.4)--(1.6) are easy and does not give them. The arguments above are the ones used in the
formalization.}}
